% !TEX root = ../b-rg.tex
\chapter{Morphophonology}\label{ch:morphophonology}
\stub

% --- Planned sections ---
% Gathers the alternations that recur across inflection and derivation, so
% the Adjectives, Verb Morphology and Word Formation chapters can
% cross-reference them instead of restating them.
%
% \section{Coda neutralisation and its reversal}
%   ð~θ, s~z, ʒ~x, v~f or vowel change (nað-ar naθ, cav-ir caf, bav-ar bo);
%   currently exemplified under Phonotactics -- move here or cross-reference.
% \section{Stem-final adjustments before vowel-initial suffixes}
%   ç -> c, c -> ch, q -> qu, consonant doubling to keep a vowel checked,
%   revoicing (tart tardessem, braf bravessem), -Vu -> -Vv-, -ous -> -os-,
%   -e -> -ey- (currently stated for comparatives; do verbs behave alike?).
% \section{Syllabic consonants in inflection}
%   -abr -> -abl-, -r -> -er- (stabr stablessem, nuvr nuverrem);
%   -mm/-nn/-rr endings (yemm, grinn, carr).
% \section{Stress-shifting suffixes}
%   -essem/-errem, -(e)nt, -(e)nç (cf. Stress); stress in compounds.
% \section{Insertion and deletion}
%   Linking -u- after -q (picquessem); epenthesis; any reduplication?
% \section{Vowel alternations}
%   Free vs checked vowels under suffixation; kw ~ k (voc voccar,
%   bocq bocquar; cf. Phonemes).
% \section{Sandhi and elision}
%   Elision (l', n', d', m', j'); article + preposition contractions
%   (dy, ag, ny; dell', all', nell', coll', sull', pall'); linking [ʀ]
%   (cf. Non-rhoticity).
